
i)
Comme @@M0@@ peut se résoudre localement par un complexe de Koszul, l'algèbre graduée @@M1@@ est alternée (2.2). Il existe donc un homomorphisme d'algèbres et un seul:
@@M2@@
prolongeant l'injection de @@M3@@ dans
@@M4@@. Il suffit de vérifier localement que c'est un isomorphisme.

On peut donc se placer sur un ouvert affine @@M5@@ de @@M6@@, @@M7@@ étant défini par un idéal @@M8@@ de @@M9@@, et se donner un homomorphisme régulier @@M10@@, définissant @@M11@@, avec @@M12@@. Soient @@M13@@ les images d'une base de @@M14@@.

Dans @@M15@@, nous noterons, pour @@M16@@, @@M17@@ l'élément @@M18@@ de @@M19@@ considéré comme étant de degré @@M20@@. On peut se restreindre à l'ouvert @@M21@@ de @@M22@@ où l'élément @@M23@@ par exemple est inversible.

Les @@M24@@ sont les faisceaux d'homologie du complexe @@M25@@. Pour étudier ce complexe sur @@M26@@, on fait le changement de base suivant dans @@M27@@: si @@M28@@ est la base donnée de @@M29@@, on pose:
@@M30@@
Posons @@M31@@; la différentielle du complexe @@M32@@ est donnée par
@@M33@@
En prenant pour base de @@M34@@ les éléments
@@M35@@ avec @@M36@@ on trouve donc:
@@M37@@
